Donald Knuth gives a formula for sums of odd powers in terms of the central factorial numbers of the second kind $T(n,k)$ in~\cite{knuth1993johann}.
\begin{proposition}[Knuth's odd powers sum]
  \label{prop:knuth-odd-powers-sum}
  Let $n$ be a nonnegative integer, and let $m$ be a positive integer.
  \begin{align*}
    \KnuthRFoldSum{1}{n}{2m-1} = \sum_{k=1}^{m} (2k-1)! \binom{n+k}{2k} T(2m,2k),
  \end{align*}
  where $T(2m,2k)$ are central factorial numbers of the second kind.
\end{proposition}
The sequences \href{https://oeis.org/A008957}{\texttt{A008957}}, \href{https://oeis.org/A036969}{\texttt{A036969}}, and \href{https://oeis.org/A269945}{\texttt{A269945}} list central factorial numbers of the second kind $T(n,k)$ in the OEIS, as described in~\cite{sloane2003line}.
Table~\ref{tab:central-factorial-numbers-t0-even} lists initial values of $T(2n,2k)$.
\begin{table}[H]
  \begin{center}
    \setlength\extrarowheight{-6pt}
    \begin{tabular}{c|cccccccc}
      $n/k$ & 0 & 1 & 2 & 3 & 4 & 5 & 6 & 7 \\
      \hline
      0 & 1 &   &      &       &       &      &    &     \\
      1 & 0 & 1 &      &       &       &      &    &     \\
      2 & 0 & 1 & 1    &       &       &      &    &     \\
      3 & 0 & 1 & 5    & 1     &       &      &    &     \\
      4 & 0 & 1 & 21   & 14    & 1     &      &    &     \\
      5 & 0 & 1 & 85   & 147   & 30    & 1    &    &     \\
      6 & 0 & 1 & 341  & 1408  & 627   & 55   & 1  &     \\
      7 & 0 & 1 & 1365 & 13013 & 11440 & 2002 & 91 & 1   \\
    \end{tabular}
  \end{center}
  \caption{Central factorial numbers of the second kind $T(2n,2k)$.}
  \label{tab:central-factorial-numbers-t0-even}
\end{table}

It is interesting to notice that Proposition~\eqref{prop:knuth-odd-powers-sum}
directly follows
from Central Newton's Formula in Lemma~\eqref{lem:central-newtons-formula-for-powers},
evaluated
at zero.
For $m\geq 1$, the formula gives
\begin{align*}
  n^{2m} = \sum_{k=1}^{2m} \frac{\centralFactorial{n}{k}}{k!} \cdot \delta^{k} 0^{2m}.
\end{align*}
The term for $k=0$ vanishes because $\delta^{0} 0^{2m} = 0$ when $m\geq1$.
For positive integers $m$ and $k$, the central difference $\delta^k 0^m$ is nonzero exactly when $k\leq m$ and $m$ and $k$ have the same parity:
\begin{align*}
  \begin{cases}
    \delta^{k} 0^m \neq 0, \quad &\text{when} \quad m \equiv k \pmod{2}, \\
    \delta^{k} 0^m = 0, \quad &\text{when} \quad m \not\equiv k \pmod{2},
  \end{cases}
\end{align*}
In particular, the difference vanishes when $m$ and $k$ have opposite parity.
This parity property applies to powers evaluated at zero; for example, $\delta^{3}1^{4}=24\neq0$.
Since $2m$ is even, the differences $\delta^{k}0^{2m}$ vanish for odd $k$, so we may omit those terms:
\begin{align*}
  n^{2m} = \sum_{k=1}^{m} \frac{\centralFactorial{n}{2k}}{(2k)!} \cdot \delta^{2k} 0^{2m}.
\end{align*}
By the definition of the central factorial numbers of the second kind, we obtain
\begin{align*}
  n^{2m} = \sum_{k=1}^{m} \centralFactorial{n}{2k} T(2m,2k).
\end{align*}
Therefore,
\begin{align*}
  n^{2m} = \sum_{k=1}^{m} (2k)! \frac{\centralFactorial{n}{2k}}{(2k)!} T(2m,2k).
\end{align*}
Applying Lemma~\eqref{lem:central-factorials-binomial-form} to simplify $\frac{\centralFactorial{n}{2k}}{(2k)!}$ gives
\begin{align*}
  n^{2m} = \sum_{k=1}^{m} (2k)!  \frac{n}{2k} \binom{n+k-1}{2k-1} T(2m,2k).
\end{align*}
Consequently,
\begin{align*}
  n^{2m} = \sum_{k=1}^{m} (2k-1)!  n \binom{n+k-1}{2k-1} T(2m,2k).
\end{align*}
Factoring out $n$ gives
\begin{align*}
  n^{2m-1}
  &= \sum_{k=1}^{m} (2k-1)! \binom{n+k-1}{2k-1} T(2m,2k).
\end{align*}
Applying the hockey-stick identity $\sum_{k=0}^{n} \binom{k}{j} = \binom{n+1}{j+1}$ gives
\begin{align*}
  \KnuthRFoldSum{1}{n}{2m-1} = \sum_{k=1}^{m} (2k-1)! \binom{n+k}{2k} T(2m,2k).
\end{align*}
This formula agrees with Proposition~\eqref{prop:knuth-odd-powers-sum}.
Repeated application of the hockey-stick identity $\sum_{k=0}^{n} \binom{k}{j} = \binom{n+1}{j+1}$ generalizes the formula to multifold sums.
\begin{proposition}[Knuth's odd powers sum multifold]
  \label{prop:knuth-odd-powers-sum-multifold}
  Let $r$ and $n$ be nonnegative integers, and let $m$ be a positive integer.
  \begin{align*}
    \KnuthRFoldSum{r+1}{n}{2m-1}
    = \sum_{k=1}^{m} (2k-1)! \binom{n+k+r}{2k+r} T(2m,2k),
  \end{align*}
  where $T(2m,2k)$ are central factorial numbers of the second kind.
\end{proposition}
The following example illustrates Proposition~\eqref{prop:knuth-odd-powers-sum}.
\begin{example}
  For non-negative integers $n$, and $r=1$, we have
  \begin{align*}
    \KnuthRFoldSum{1}{n}{1} &= \tbinom{n+1}{2} 1, \\
    \KnuthRFoldSum{1}{n}{3} &= \tbinom{n+2}{4} 6    + \tbinom{n+1}{2} 1, \\
    \KnuthRFoldSum{1}{n}{5} &= \tbinom{n+3}{6} 120  + \tbinom{n+2}{4} 30   + \tbinom{n+1}{2} 1, \\
    \KnuthRFoldSum{1}{n}{7} &= \tbinom{n+4}{8} 5040 + \tbinom{n+3}{6} 1680 + \tbinom{n+2}{4} 126 + \tbinom{n+1}{2} 1.
  \end{align*}
\end{example}

The following example illustrates Proposition~\eqref{prop:knuth-odd-powers-sum-multifold}.
\begin{example}
  For non-negative integers $r,n$
  \begin{align*}
    \KnuthRFoldSum{r}{n}{1} &= \tbinom{n+0+r}{1+r} 1, \\
    \KnuthRFoldSum{r}{n}{3} &= \tbinom{n+1+r}{3+r} 6    + \tbinom{n+0+r}{1+r} 1, \\
    \KnuthRFoldSum{r}{n}{5} &= \tbinom{n+2+r}{5+r} 120  + \tbinom{n+1+r}{3+r} 30   + \tbinom{n+0+r}{1+r} 1, \\
    \KnuthRFoldSum{r}{n}{7} &= \tbinom{n+3+r}{7+r} 5040 + \tbinom{n+3+r}{5+r} 1680 + \tbinom{n+1+r}{3+r} 126 + \tbinom{n+0+r}{1+r} 1.
  \end{align*}
\end{example}

The coefficients in Proposition~\eqref{prop:knuth-odd-powers-sum-multifold} form the sequence \href{https://oeis.org/A303675}{\texttt{A303675}} in the OEIS, as listed in~\cite{sloane2003line}.
Proposition~\eqref{prop:knuth-odd-powers-sum-multifold} also motivates the following identity.
Combining Newton's formula in Lemma~\eqref{lem:central-newtons-formula-for-powers} with Lemma~\eqref{lem:central-factorials-binomial-form} gives
\begin{align*}
  n^m = \sum_{k=1}^{m} \frac{n}{k} \binom{n+\frac{k}{2}-1}{k-1} \delta^{k} 0^m.
\end{align*}
Factoring out $n$ gives
\begin{align*}
  n^{m-1} = \sum_{k=1}^{m} \frac{1}{k} \binom{n+\frac{k}{2}-1}{k-1} \delta^{k} 0^m.
\end{align*}
Shifting the exponent gives
\begin{align*}
  n^{m} = \sum_{k=1}^{m+1} \frac{1}{k} \binom{n+\frac{k}{2}-1}{k-1} \delta^{k} 0^{m+1}.
\end{align*}
Re-indexing the sum gives
\begin{align}
  \label{eq:power-identity-at-zero}
  n^{m} = \sum_{k=0}^{m} \frac{1}{k+1} \binom{n+\frac{k}{2}-\frac{1}{2}}{k} \delta^{k+1} 0^{m+1}.
\end{align}
Setting $r=0$ and $t=\frac{1}{2}$ in Lemma~\eqref{lem:centered-hockey-stick-identity} gives the following proposition.
\begin{proposition}[Ordinary sums of powers at zero]
  \label{prop:ordinary-sums-of-powers-at-zero}
  Let $n$ and $m$ be nonnegative integers.
  \begin{align*}
    \KnuthRFoldSum{1}{n}{m}
    = \sum_{k=0}^{m}
    \left[ \binom{n+\frac{k}{2}+\frac{1}{2}}{k+1} - \binom{\frac{k}{2}+\frac{1}{2}}{k+1}\right]
    \frac{\delta^{k+1} 0^{m+1}}{k+1}.
  \end{align*}
\end{proposition}
The formula can also be written in terms of central factorial numbers of the second kind:
\begin{align*}
  \KnuthRFoldSum{1}{n}{m}
  = \sum_{k=0}^{m} k!
  \left[ \binom{n+\frac{k}{2}+\frac{1}{2}}{k+1} - \binom{\frac{k}{2}+\frac{1}{2}}{k+1}\right]
  T(k+1, m+1).
\end{align*}
Setting $r=1$ and $t=\frac{1}{2}$ in Lemma~\eqref{lem:centered-hockey-stick-identity} gives a formula for double sums of powers.
\begin{proposition}[Double sums of powers at zero]
  \label{prop:double-sums-of-powers-at-zero}
  Let $n$ and $m$ be nonnegative integers.
  \begin{align*}
    \KnuthRFoldSum{2}{n}{m}
    = \sum_{k=0}^{m}
    \left[
      \binom{n+\frac{k}{2}+\frac{3}{2}}{k+2}
      - \binom{\frac{k}{2}+\frac{1}{2}}{k+1} \KnuthRFoldSum{1}{n}{0}
      - \binom{\frac{k}{2}+\frac{3}{2}}{k+2} \KnuthRFoldSum{0}{n}{0}
    \right]
    \frac{\delta^{k+1} 0^{m+1}}{k+1}.
  \end{align*}
\end{proposition}
The general formula is as follows.
\begin{theorem}[Multifold sums of powers at zero]
  \label{thm:multifold-sums-of-powers-at-zero}
  Let $r$, $n$, and $m$ be nonnegative integers.
  \begin{align*}
    \KnuthRFoldSum{r}{n}{m}
    =
    \sum_{k=0}^{m}
    \left[
      \binom{n+\frac{k}{2}+\frac{2r-1}{2}}{k+r}
      -
      \sum_{s=0}^{r-1} \binom{\frac{k}{2} + \frac{2(r-s)-1}{2}}{k+r-s}
      \KnuthRFoldSum{s}{n}{0}
    \right] \frac{\delta^{k+1} 0^{m+1}}{k+1}.
  \end{align*}
  \begin{proof}
    By repeatedly applied Lemma~\eqref{lem:centered-hockey-stick-identity},
    and Equation~\eqref{eq:power-identity-at-zero}.
  \end{proof}
\end{theorem}
Combining Lemma~\eqref{lem:multifold-sum-of-zero-powers} with Theorem~\eqref{thm:multifold-sums-of-powers-at-zero} gives the following binomial form.
\begin{proposition}[Multifold sums of powers at zero binomial form]
  \label{prop:multifold-sums-of-powers-at-zero-binomial-form}
  Let $r$, $n$, and $m$ be nonnegative integers.
  \begin{align*}
    \KnuthRFoldSum{r}{n}{m}
    =
    \sum_{k=0}^{m}
    \left[
      \binom{n+\frac{k}{2}+\frac{2r-1}{2}}{k+r}
      -
      \sum_{s=0}^{r-1} \binom{\frac{k}{2} + \frac{2(r-s)-1}{2}}{k+r-s}
      \binom{s+n-1}{s}
    \right] \frac{\delta^{k+1} 0^{m+1}}{k+1}.
  \end{align*}
  \begin{proof}
    By repeatedly applied Lemma~\eqref{lem:centered-hockey-stick-identity},
    and Equation~\eqref{eq:power-identity-at-zero}.
  \end{proof}
\end{proposition}
This formula can also be expressed in terms of central factorial numbers of the second kind:
\begin{align*}
  \KnuthRFoldSum{r}{n}{m}
  =
  \sum_{k=0}^{m}
  \left[
    \binom{n+\frac{k}{2}+\frac{2r-1}{2}}{k+r}
    -
    \sum_{s=0}^{r-1} \binom{\frac{k}{2} + \frac{2(r-s)-1}{2}}{k+r-s}
    \binom{s+n-1}{s}
  \right] k! T(k+1, m+1).
\end{align*}
The results above follow from evaluating the central Newton formula at zero.
Table~\ref{tab:central-factorial-numbers-t0} lists initial values of $T(n,k)$.
\begin{table}[H]
  \begin{center}
    \setlength\extrarowheight{-6pt}
    \begin{tabular}{c|ccccccccc}
      $n/k$ & 0 & 1 & 2 & 3 & 4 & 5 & 6 & 7 & 8 \\ [3px]
      \hline
      0 & 1 &   &   &   &   &   &   &   &   \\
      1 & 0 & 1 &   &   &   &   &   &   &   \\
      2 & 0 & 0 & 1 &   &   &   &   &   &   \\
      3 & 0 & $\tfrac{1}{4}$ & 0 & 1 &   &   &   &   &   \\
      4 & 0 & 0 & 1 & 0 & 1 &   &   &   &   \\
      5 & 0 & $\tfrac{1}{16}$ & 0 & $\tfrac{5}{2}$ & 0 & 1 &   &   &   \\
      6 & 0 & 0 & 1 & 0 & 5 & 0 & 1 &   &   \\
      7 & 0 & $\tfrac{1}{64}$ & 0 & $\tfrac{91}{16}$ & 0 & $\tfrac{35}{4}$ & 0 & 1 &   \\
      8 & 0 & 0 & 1 & 0 & 21 & 0 & 14 & 0 & 1 \\
    \end{tabular}
  \end{center}
  \caption{Central factorial numbers of the second kind $T(n,k)$.
    See also the sequences \href{https://oeis.org/A395862}{\texttt{A395862}},
    and \href{https://oeis.org/A370703}{\texttt{A370703}} in the OEIS~\cite{sloane2003line}.
  }
  \label{tab:central-factorial-numbers-t0}
\end{table}

